Agents connected to many tool servers retrieve a few candidate tools and let a language model
pick one. Evaluations of this step count every wrong tool alike, but real catalogs pair each read
operation with tools that create, send, overwrite or delete on the same resource, and choosing one
of those when the user only asked to look would be an incident, not a near miss. We measure this
selection-time \emph{side-effect escalation} on the public Model Context Protocol ecosystem:
\NumRawTools{} tool definitions harvested from \NumRespondingServers{} registry servers, a
\NumCatTools{}-tool catalog with \NumSibGroups{} sibling groups, and a targeted benchmark of
\NumQueries{} validated requests. A tool's nearest neighbour on its own server in embedding space is a sibling
\NNSibQwen\,\% of the time (\NNChance\,\% by chance). Routing to the top retrieved tool selects a tool that changes state for
\RWRetRate\,\% of read requests at \NumPoolTools{} tools; LLM selection over a
shortlist lowers this to \RWQPlainRate\,\% for Qwen3.5, and overall escalation from
\GRNonePoolConf\,\% to \SelQPlainPoolConf\,\% of requests, but never removes it. Automatic severity
labels are unreliable: only \EscPrec\,\% of the escalations they flag survive human review, most often
because servers mark ordinary overwrites as destructive. Whether a pick ``does more than asked'' is
itself contested: the two authors, rating independently, agree only slightly ($\kappa=\IndKappa$), so we
report human-reviewed labels, with independent double-rating and consensus on a stratified subset, and show
how the estimates move with the rater. Side-effect annotations in
the prompt barely change what the selector picks; adding each candidate's sibling group raises
accuracy by 2.6--4.5 points for the larger selectors, and for Gemma3 and Phi-4 a plain shortlist of the same length does not
come close, so a longer list alone does not explain the gain; an escalation guard removes most remaining
escalation at the price of holding 15--20\,\% of calls for confirmation. We derive guidance for agent builders
on routing, annotations and confirmation policies.
